\documentclass[12pt]{article}
\usepackage{samstyle} 

\newcommand{\NuevaPagina}[1]{\newpage
\mbox{
\begin{minipage}{3cm}
    \includegraphics[width=3cm]{logo}
\end{minipage}
\begin{minipage}{14.7cm}
  \begin{flushright}
\hfill P�gina #1 \hfill Primer parcial-ST0240 \hfill Marzo
\hfill C�digo: \rule{4cm}{0.3pt}
  \end{flushright}
\end{minipage}}\vspace{3pt}

\rule{18cm}{1.5pt}}

%%% COMIENZO DE DOCUMENTO %%%
\begin{document}
\pagestyle{empty}
%%% Encabezado
\mbox{
    \begin{minipage}{5cm}
        \includegraphics[width=5cm]{logo}
    \end{minipage}
    \begin{minipage}{13cm}
    \begin{center}
        \textsl{
        \textbf{\Large Escuela de Ingenier�as}\\
        \textbf{\large Programaci�n de Computadores - ST0240}\\
        \textbf{\large Pr�ctica Final (Valor 20\%)}}
    \end{center}
    \end{minipage}
}
\begin{flushright}
  \begin{minipage}{4.5cm}
    Nota: \hrulefill
  \end{minipage}
\end{flushright}\vspace{-2mm}

% Nombre: \rule{9cm}{0.3pt} \; C�digo: \hrulefill\vspace{6pt}

Profesor: \underbar{\hspace{1cm} Sergio Monsalve
\hspace{3.1cm}} \; \hfill Grupo: \underbar{\hspace{1cm} 063 } \;
Mayo \underbar{\hspace{1cm}} de 2016 \rule{18cm}{1.5pt}
%%% Fin del encabezado


La siguiente es la descripci�n de la Practica Final del curso de Programaci�n de computadores. La fecha de entrega es el d�a domingo 5 de Junio a las 11:59pm a trav�s de la plataforma de EAFIT Interactiva.\\

Los empleados de cierta empresas registran la entrada y salida del trabajo de manera electr�nica, cada empleado tiene un c�digo para poder ser identificado y saber el n�mero de horas que ha trabajado. La empresa tiene pol�ticas de compensaci�n para los empleados que trabajan horas extras, y para ello tiene horarios de trabajo.\\

La empresa desea saber cual es la cantidad de horas y el promedio de horas que los empleados est�n trabajando al mes, tambien el n�mero de horas extras que est�n trabajando los empleados y las horas nocturnas. Para ello tiene las siguientes pol�ticas en cuanto a los horarios:\\

La entrada de Lunes a S�bado es despu�s de las 6 am.\\

Horas Extra: \\
\begin{itemize}
  \item  Si alguien trabaja despu�s de las 4 pm los lunes se cuentan estas como horas extras, 
  \item  De martes a viernes despu�s de las 3 pm se contar�n como extras. 
  \item  Los s�bados estas ser�n contadas despu�s de la 1:45 pm. 
  \item  Los domingos son especiales, ya que si se trabaja ese d�a, las horas laboradas son todas horas extras. 
\end{itemize}

Horas Nocturnas: \\

\begin{itemize}
  \item Tambi�n se cuenta con otro de horario especial, las horas laboradas entre las 10 pm y las 6 am ser�n consideradas como nocturnas.
\end{itemize}

\section*{Entrada}

Una muestra del formato del archivo de entrada es el siguiente:

\begin{itemize}
  \item Los valores est�n divididos por comas(",") y por fin de linea("{\textbackslash}n"):
  \item El primer valor, puede ser 1, 2 o 3, si es 1 es la hora de entrada y los valores 2 y 3 marcan la hora de salida.
  \item El siguiente par�metro es el c�digo del empleado.
  \item El siguiente par�metro es el �rea donde se trabaja dentro de la empresa.
  \item El siguiente par�metro es la hora y el septimo es el minuto.
  \item El siguiente valor es el mes, luego es el d�a y el ultimo es el a�o.
\end{itemize}

\section*{Salida}

Se requiere que usted elabore un programa en VBA para excel, que cumpla con los siguientes requisitos: 

\begin{enumerate}
\item \textbf{(Valor 25\%)} Calcule las horas extra trabajadas:
\end{enumerate}

\begin{itemize}
  \item Cantidad de horas extras por empleado
  \item Cantidad de horas extras trabajadas de todos los empleados
  \item Promedio de horas extras trabajadas por los empleados
\end{itemize}
\vspace{0.5cm}

\begin{enumerate}[resume] 
\item \textbf{(Valor 25\%)} Calcule las horas nocturnas trabajadas:
\end{enumerate}

\begin{itemize}
  \item Cantidad de horas nocturnas por empleado
  \item Cantidad de horas nocturnas trabajadas de todos los empleados
  \item Promedio de horas nocturnas trabajadas por los empleados
\end{itemize}

\vspace{0.5cm}

\begin{enumerate}[resume]
\item \textbf{(Valor 25\%)} Calcule las horas totales trabajadas:
\end{enumerate}

\begin{itemize}
  \item Cantidad de horas trabajadas por empleado
  \item Cantidad de horas trabajadas de todos los empleados
  \item Promedio de horas trabajadas por los empleados
\end{itemize}

\vspace{0.5cm}

\begin{enumerate}[resume]
\item \textbf{(Valor 25\%)} Visualizaci�n, interpretaci�n y reportes: genere las siguientes gr�ficas:
\end{enumerate}

\begin{itemize}
  \item  Una gr�fica para las horas extras 
  \item  Una gr�fica para las horas nocturnas 
  \item  Una gr�fica para el total de horas trabajadas
\end{itemize}

\end{document}
